\chapter{Introducción}
\label{cap:introduccion}

\begin{definicion}
  Un vector del espacio se escribe en la base canónica como
  $\vec{v} = a\,\ihat + b\,\jhat + c\,\khat$.
\end{definicion}

\begin{teorema}[Norma euclidiana]
  \label{teo:norma}
  Para todo $\vec{v} = a\,\ihat + b\,\jhat + c\,\khat$ se cumple
  \[
    \lVert \vec{v} \rVert = \sqrt{a^2 + b^2 + c^2}.
  \]
\end{teorema}

\begin{proof}
  Aplicar dos veces el teorema de Pitágoras.
\end{proof}

\begin{ejemplo}
  Si $\vec{v} = 3\,\ihat + 4\,\jhat$, entonces $\lVert \vec{v} \rVert = 5$
  (teorema~\ref{teo:norma}).
\end{ejemplo}

\begin{observacion}
  Agrega más capítulos en la carpeta \texttt{capitulos/} e inclúyelos con \verb|\input|.
\end{observacion}
